\documentclass[12pt]{article}
\usepackage[spanish]{babel}
\usepackage{amsmath,amssymb}
\usepackage{geometry}
\geometry{margin=2.5cm}

\begin{document}

\section*{I. Halle la fórmula para cada una de las siguientes sumatorias}

\begin{enumerate}
\item $\displaystyle \sum_{i=1}^{n} i^4$
\item $\displaystyle \sum_{i=1}^{n} i^5$
\item $\displaystyle \sum_{i=1}^{n} \left(\sqrt{9+x}\right)^i$
\item $\displaystyle \sum_{i=1}^{n} \frac{1}{16i^2-1}$
\item $\displaystyle \sum_{i=1}^{n} \tan(aix)$
\item $\displaystyle \sum_{i=1}^{n} sen^i(ax)$
\end{enumerate}

---

\section*{II. Encontrar el área de la región exacta, expresar como límite de sumas de Riemann, con particiones iguales}

\begin{enumerate}
\setcounter{enumi}{6}
\item Hallar el área de la región acotada por $y=5x-6$, el eje $x$ y las rectas $x=1$ y $x=5$.
\item Hallar el área de la región acotada por $y=5\sqrt{x}$, el eje $x$ y las rectas $x=1$ y $x=4$.
\item Hallar el área de la región acotada por $f(x)=16-x^2$, el eje $x$ y las rectas $x=1$ y $x=4$.
\item Hallar el área de la región acotada por $y=5-|x|$, el eje $x$ y las rectas $x=-5$ y $x=5$.
\item Hallar el área de la región acotada por $y=\tan x$, el eje $x$ y las rectas $x=0$ y $x=\frac{\pi}{3}$.
\end{enumerate}

---

\section*{III. Halle las siguientes integrales definidas}

\begin{enumerate}
\setcounter{enumi}{11}
\item $\displaystyle \int_{1}^{3}\left(\frac{1}{2x-1}-\frac{1}{2x+1}\right)\,dx$
\item $\displaystyle \int_{2}^{4}\frac{e^{\ln x}}{x^2+7}\,dx$
\item $\displaystyle \int_{0}^{1}\frac{arctg x}{x^2+1}\,dx$
\item $\displaystyle \int_{6}^{9}\frac{3\ln x-5}{x}\,dx$
\end{enumerate}

---

\section*{IV. Calcule el área de la región acotada por las curvas dadas}

\begin{enumerate}
\setcounter{enumi}{15}
\item $y=x^3$, $y=0$, $x=-1$, $x=3$
\item $y=4x-x^2$, $y=0$, $x=1$, $x=3$
\item $x=1+y^2$, $x=10$
\item $x=y^2+4y$, $x=0$
\item Un sector circular de radio $r$ y ángulo $\alpha$
\item Dentro de $y=25-x^2$, $256x=3y^2$, $16y=9x^2$
\item $y=x^4-4x^2$, $y=4x^2$
\item $y=e^{x/a}+e^{-x/a}$, $y=0$, $x=\pm a$
\item $y=\frac{1}{1+x^2}$, $y=0$, $x=\pm1$
\item \item Encontrar el área comprendida entre las curvas $y=x^2$ y $y=x^3$.
\item Encontrar el área limitada por la curva $y^2-2y+x-8$, las rectas $x=0$, $y=-1$, $y=3$.
\item Hallar el área comprendida entre las curvas $y=2x$, $y=\frac{1}{2}(x^3)$ en el primer cuadrante, usando rectángulos diferenciales horizontales.
\item Usando integrales, calcular el área de la elipse con centro en el origen, semieje mayor $a$ y semieje menor $b$.
\end{enumerate}

---

\section*{V. Encontrar el volumen generado en la rotación del área plana dada alrededor del eje indicado}

\begin{enumerate}
\setcounter{enumi}{28}

\item $y=x^3$, $y=0$, $x=2;x=2$.
\item $y^2=x^3$, $y=0$, $x=2$; eje $x$.
\item Calcule el volumen de un cono de radio $r$ y altura $h$.
\item Calcule el volumen generado por una elipse horizontal con semieje mayor $a$ y semieje menor $b$ al rotar alrededor del eje $x$.
\item Hallar el volumen que se genera al rotar el área comprendida en la parábola $y=4x-x^2$ y el eje $x$ con respecto a la recta $y=6$.
\item Encontrar el volumen del sólido que se genera al girar la región comprendida entre $y=e^x$, $y=1$, $x=1$ alrededor del eje $x$.
\item Hallar el volumen del toroide que se genera al girar la curva $(x-3)^2+y^2=4$ alrededor del eje $x$.
\item Hallar el volumen del sólido que se genera al girar la región comprendida entre las curvas $y=8x^2$ y $y=\sqrt{8x}$ alrededor del eje $x$.
\item Use el método más apropiado para calcular el volumen que se genera al girar el área limitada por las curvas $y=2-x^4$, $y=1$ alrededor del eje $x$.
\end{enumerate}

---

\section*{VI. Calcule las siguientes integrales impropias}

\begin{enumerate}
\setcounter{enumi}{37}
\item $\displaystyle \int_{1}^{+\infty} \frac{dx}{x^2}$
\item $\displaystyle \int_{-\infty}^{0} \frac{dx}{(4-x)^2}$
\item $\displaystyle \int_{0}^{1} \frac{dx}{\sqrt{x}}$
\item $\displaystyle \int_{-2}^{2} \frac{dx}{\sqrt{4-x^2}}$
\item $\displaystyle \int_{0}^{4} \frac{dx}{4-x}$
\item $\displaystyle \int_{-\infty}^{0} xe^x\,dx$
\item $\displaystyle \int_{-\infty}^{+\infty} \frac{dx}{1+4x^2}$
\end{enumerate}

---

\section*{VII. Halle la longitud de curva o del arco indicado}

\begin{enumerate}
\setcounter{enumi}{44}
\item $y^2=8x^2$ desde $x=1$ a $x=8$.
\item $27y^2=4(x-2)^3$ desde $(2,0)$ a $(11,6\sqrt{3})$.
\item $y=\ln x$ desde $x=1$ a $x=2\sqrt{2}$.
\item $y=\ln(e^x-1)/(e^x+1)$ desde $x=2$ a $x=4$.
\item $y=\ln(\cos x)$ desde $x=\pi/6$ a $x=\frac{1}{4} \pi$.
\item Calcular la longitud del arco de circunferencia de radio $3$ que forma un ángulo de $60^\circ$ con el eje $x$ positivo. Comprobar la respuesta usando la fórmula para cálculo del arco de circunferencia conociendo ángulo y radio.
\end{enumerate}

\end{document}